\section{符号说明}\label{sec:notation}

本文符号与 \S\ref{sec:analysis} 式~\eqref{eq:pipeline} 一致，主要符号如表~\ref{tab:notation}。

\begin{table}[H]
\centering
\caption{主要符号与含义。}
\label{tab:notation}
\footnotesize
\setlength{\tabcolsep}{4pt}
\resizebox{\linewidth}{!}{\begin{tabular}{@{}lll@{}}
  \toprule
  符号 & 含义 & 单位 / 取值 \\
  \midrule
  $G$ & 二倍体基因组（碱基串与调控基序） & 512\,bp \\
  $q(G)$ & 基因组经调控网络的读出 & $[0,1]^{8}$ \\
  $\mathbf g_i$ & 第 $i$ 个细胞的调控网络状态 & $\mathbb{R}^{8}$ \\
  $M_k$ & 第 $k$ 个调控基序 & 6\,bp，共 8 个 \\
  $D_{\Theta}$ & 发育算子（调控网络、细胞命运、连接生成） & --- \\
  $N$ & 活跃神经元数 & 整数，$[24,48]$ \\
  $\mathbf p_i$ & 第 $i$ 个细胞在发育方域内的位置 & 归一化单位 \\
  $A$ & 连接邻接矩阵（二值） & $N\times N$ \\
  $Z$ & 细胞身份（六类功能域的后验） & 概率向量 \\
  $\tau_i$ & 第 $i$ 个神经元的整合时间常数 & $[1,10]$ \\
  $W^{(0)}$ & 发育初始连接权重 & 实矩阵 \\
  $M$ & 活跃神经元掩码 & 二值向量 \\
  $h_i^{t}$ & 第 $i$ 个神经元在 $t$ 步的状态 & $(-1,1)$ \\
  $x_t$ & 第 $t$ 步的观测向量 & $[0,1]^{12}$ \\
  $U_i,\,m_i,\,b_i$ & 感官增益、输入权重、偏置（按细胞类型固定） & 实数 \\
  $\phi$ & 神经元非线性 & $\tanh$ \\
  $\omega_t,\,v_t$ & 第 $t$ 步的转向与推进动作 & $\omega\in[-1,1],\ v\in[0,1]$ \\
  $\Theta$ & 当代学习优化的参数张量 & 与 $W^{(0)}$ 同形 \\
  $W$ & 实际连接权重 & $W=A\odot(\operatorname{sign}(W^{(0)})\odot\mathrm{softplus}(\Theta))$ \\
  $\Delta W$ & 当代学习产生的权重增量（不遗传） & 实矩阵 \\
  $L$ & 当代学习算子 & --- \\
  $F$ & 适应度 & 代内 $[0,1]$ \\
  $S,P,E_{\mathrm{esc}},Q$ & 生存、捕食、逃脱、净能量四分量 & $[0,1]$ \\
  $E_t$ & 能量 & $[0,E_{\max}]$，$E_{\max}=1$ \\
  $H_t$ & 饥饿 & $[0,1]$，$H_t=1-E_t/E_{\max}$ \\
  $\lambda,\,\gamma$ & 空间布线代价、表达相似度系数 & $\lambda=2.0$，$\gamma=1.0$ \\
  $\rho$ & 调控网络状态更新率 & $0.35$ \\
  $\mu$ & 每碱基每配子的突变率 & $0.001$ \\
  $T$ & 回合长度 & $600$ 步（$20$\,Hz，$30$\,s） \\
  \bottomrule
\end{tabular}}
\end{table}
